\documentclass[12pt]{article}
\usepackage[utf8]{inputenc}
\usepackage[T1]{fontenc}
\usepackage{hyperref}
\title{Silt'e Language Resources for K-12 Curriculum Development}
\author{Compiled by Claude Sonnet 5.5 (claude-sonnet-5-5) for Andr\'as Kornai}
\date{}
\begin{document}
\maketitle

\section{Introduction}
Silt'e (ISO 639-3 \texttt{stv}, Glottocode silt1240) is a Semitic language of Ethiopia, spoken mainly in the Silt'e Zone. Sources in search results give about 880,000 speakers (2007 census), far above the figure of about 1,000 in the task description. It is written in the Ge'ez script. Web search found mostly academic studies and little in the way of textbooks.

\section{Primary Educational Resources}
No specific Silt'e textbooks or primers were found in search results.

\section{Linguistic and Reference Materials}
\begin{itemize}
\item Wikipedia, Silt'e language: \url{https://en.wikipedia.org/wiki/Silt%CA%BCe_language}
\item Wikipedia, Silt'e people: \url{https://en.wikipedia.org/wiki/Silt%CA%BCe_people}
\end{itemize}

\section{Religious and Community Resources}
None found.

\section{Digital and Audio Resources}
None found.

\section{Curriculum Development Context}
\begin{itemize}
\item Addis Ababa University study, ``Language Planning and Policy in the Silt'e Zone'' (2015). It finds Silt'e is used as a medium of instruction, but is weak in personal and public communication, and speakers feel they need Amharic: \url{https://etd.aau.edu.et/items/659cdb9c-1c4a-4d10-af3c-a93649760a1f}
\item Study on mother-tongue medium of instruction in Ethiopian primary schools. It identifies textbook shortages as a major obstacle, but covers Ethiopia generally and not Silt'e alone: \url{https://glocal.soas.ac.uk/afala2023-121/}
\end{itemize}
Textbook supply is the main gap, and the actual textbook holdings are not documented in what I found.

\end{document}